% Auto-generated by final-edits/evaluate_with_mcc.py
% Main NER table with MCC (entity-detection binary MCC recommended for imbalanced reporting)

% --- Table: tab:main_ner (replace in paper) ---
\begin{table}[t]
\centering
\caption{Main NER performance on validation and test splits (entity-level P/R/F1; macro-F1; token accuracy; MCC).}
\label{tab:main_ner}
\footnotesize
\setlength{\tabcolsep}{3pt}
\resizebox{\columnwidth}{!}{%
\begin{tabular}{@{}lcccccc@{}}
\toprule
\textbf{Split} & \textbf{Prec.} & \textbf{Rec.} & \textbf{F1} & \textbf{Macro-F1} & \textbf{Acc.} & \textbf{MCC} \\
\midrule
Validation & 0.702 & 0.714 & 0.708 & 0.682 & 0.954 & 0.774 \\
Test         & 0.698 & 0.703 & 0.701 & 0.677 & 0.952 & 0.776 \\
\bottomrule
\end{tabular}%
}
\vspace{0.25em}

\noindent{\scriptsize \textbf{Notes:} F1 is entity-level (seqeval strict IOB2). Acc.\ is token-level. MCC is binary entity-detection (\texttt{O} vs.\ any entity token), computed from the confusion matrix below; token-level 31-class MCC: validation=0.731, test=0.734.}
\end{table}

% --- Table: tab:confusion_binary (new quantitative confusion matrix) ---
\begin{table}[t]
\centering
\footnotesize
\caption{Token-level confusion matrix for entity detection (\texttt{O} vs.\ entity) on the test split.}
\label{tab:confusion_binary}
\begin{tabular}{@{}lrr@{}}
\toprule
 & \textbf{Predicted O} & \textbf{Predicted Entity} \\
\midrule
True O & 102,642 & 1,842 \\
True Entity & 2,630 & 8,783 \\
\bottomrule
\end{tabular}
\end{table}
